\chapter{Đánh giá, hạn chế và hướng phát triển}
\label{c:evaluation}

Chương này tổng hợp kết quả của quá trình phân tích và thiết kế hệ thống OISM,
đồng thời trình bày các hạn chế và hướng phát triển trong tương lai.

\section{Kết quả đạt được}

Hệ thống được định hướng giải quyết bài toán quản lý bán hàng và tồn kho trong môi trường
đa kênh. Các nội dung chính đã được phân tích và thiết kế gồm quản lý sản phẩm,
quản lý tồn kho, quản lý đơn hàng, cơ chế giữ hàng, POS và báo cáo.

\section{Mức độ đáp ứng yêu cầu}

Các yêu cầu chức năng và phi chức năng được trình bày trong Chương 2 và được liên kết
với các nội dung thiết kế ở các chương tiếp theo. Cách tổ chức này giúp nhóm dễ đối chiếu
giữa yêu cầu ban đầu và giải pháp thiết kế.

\begin{table}[htbp]
    \centering
    \caption{Tổng hợp mức độ đáp ứng yêu cầu}
    \label{tab:evaluation-requirements}
    \begin{tabular}{|p{0.25\textwidth}|p{0.30\textwidth}|p{0.30\textwidth}|}
        \hline
        \textbf{Nhóm yêu cầu} & \textbf{Nội dung chính} & \textbf{Đánh giá} \\
        \hline
        Xác thực và phân quyền & Người dùng, vai trò, tenant & Đáp ứng theo thiết kế \\
        \hline
        Sản phẩm & Sản phẩm, SKU, biến thể, mã vạch & Đáp ứng theo thiết kế \\
        \hline
        Tồn kho & Ledger, nhập, xuất, chuyển kho & Đáp ứng theo thiết kế \\
        \hline
        Đơn hàng & Đơn hàng tập trung, trạng thái & Đáp ứng theo thiết kế \\
        \hline
        Giữ hàng & Kiểm tra tồn và giữ hàng & Đáp ứng theo thiết kế \\
        \hline
        POS & Tìm kiếm, kiểm tra tồn, thanh toán & Đáp ứng theo thiết kế \\
        \hline
        Báo cáo & Doanh thu, lợi nhuận, tồn kho & Đáp ứng theo thiết kế \\
        \hline
    \end{tabular}
\end{table}

\section{Ưu điểm của giải pháp}

Một số ưu điểm chính của giải pháp gồm:

\begin{itemize}
    \item Quản lý tập trung dữ liệu bán hàng và tồn kho từ nhiều kênh.
    \item Thiết kế Inventory Ledger theo hướng chỉ ghi thêm các biến động.
    \item Có cơ chế giữ hàng nhằm hạn chế tình trạng bán vượt tồn.
    \item Hỗ trợ mô hình đa tenant phù hợp với hệ thống SaaS.
    \item Kiến trúc được tổ chức theo hướng rõ ràng, dễ mở rộng.
\end{itemize}

\section{Hạn chế}

Trong phạm vi đồ án, hệ thống vẫn có một số hạn chế:

\begin{itemize}
    \item Chưa bao phủ đầy đủ mọi nghiệp vụ đặc thù của từng sàn thương mại điện tử.
    \item Các yêu cầu về hiệu năng và tải lớn cần được kiểm chứng bằng môi trường triển khai thực tế.
    \item Một số giao diện và quy trình có thể cần tiếp tục điều chỉnh sau khi lấy ý kiến người dùng.
    \item Khả năng tích hợp với các dịch vụ bên ngoài phụ thuộc vào API và chính sách của từng nền tảng.
\end{itemize}

\section{Hướng phát triển}

Trong tương lai, hệ thống có thể được mở rộng theo các hướng:

\begin{itemize}
    \item Tích hợp thêm nhiều kênh bán hàng và đơn vị vận chuyển.
    \item Bổ sung ứng dụng di động cho người quản lý và nhân viên.
    \item Cải thiện các báo cáo và khả năng phân tích dữ liệu.
    \item Bổ sung chức năng dự báo nhu cầu và hỗ trợ quyết định nhập hàng.
    \item Mở rộng khả năng xử lý khi số lượng tenant và giao dịch tăng cao.
    \item Hoàn thiện các cơ chế giám sát, sao lưu và phục hồi hệ thống.
\end{itemize}

\section{Kết luận}

Qua quá trình phân tích và thiết kế, đề tài đã xây dựng được một mô hình hệ thống
quản lý bán hàng và tồn kho đa kênh với các thành phần và quy trình chính được xác định
rõ ràng. Các thiết kế về yêu cầu, kiến trúc, dữ liệu, nghiệp vụ, UML và giao diện tạo thành
cơ sở để tiếp tục hoàn thiện hệ thống trong các giai đoạn tiếp theo.
